\documentclass[11pt]{article}
\usepackage{amsmath,amssymb}
\usepackage{graphicx}
\usepackage[margin=1.1in]{geometry}
\usepackage{parskip}
\setlength{\parindent}{0pt}
\title{Companion A to the Blank-Slate Audit: Objections Answered, Tolerances Made Explicit,\\ and Sources Graded by Kind\\\\
\large New computations, an objection audit, and a concept ledger for ``The Universe and Consciousness''}
\author{Meriem}
\date{11 November 2026}
\begin{document}
\maketitle

\begin{abstract}
This document is a companion to \emph{The Universe and Consciousness: A Blank-Slate Audit} (43 pages,
13 sections, 18 figures). It does three things that the main paper had no room for. First, it answers
the objections a reader will actually raise: eighteen of them, each answered with either a number I
computed in this pass or an explicit statement that I cannot answer it and why. Second, it makes my
own choices measurable. Where the main paper said ``the coincidence is not sharp,'' this companion
computes the whole curve of how sharp it is as a function of the tolerance I am free to grant, and
where the main paper quoted a pooled $H_0$, this companion shows what every defensible pooling
estimator returns on the same four numbers, including a fully marginalised Bayesian model, and how far
the pooled value moves when each experiment is deleted in turn. Third, it grades the source pool by
\emph{kind of object} --- dataset, peer-reviewed primary, peer-reviewed reanalysis, preprint, curated
reference work, blog --- with a worked verification of one case (the PBR theorem) that also produced a
small lesson about withdrawn comments. Five new figures, all generated from the computations
themselves, are embedded. Everything here is reproducible with one command, named in the appendix.

\vspace{4pt}
\hrule
\vspace{4pt}
\textbf{Provenance note, stated up front because it affects what is and is not re-verified here.}
Between the second and third pass the sandbox that held my working tree was reset: the scripts, logs
and the 18 embedded PNGs of the main paper were lost from the workspace (they survive in the uploaded
evidence kit and inside the published PDF, which is the reason I uploaded a kit at all). The analysis
in this companion was therefore rebuilt from the raw data files --- re-downloaded from the same GitHub
release --- in a new script, \texttt{deepen.py}, and cross-checked against the numbers printed in the
main PDF. Section~3.4 reports what that independent reimplementation caught, in both directions.
\end{abstract}

\section{What this pass was for}

The brief I was given was to deepen the audit, diversify and grade the sources, deepen the concepts,
and leave the reader no opening for a question. That last instruction is impossible in the literal
reading and worth taking seriously in the intended one: I cannot prevent questions, but I can
\emph{pre-empt} them, which means anticipating the objection, answering it with the cheapest
computation that settles it, and writing down the answers I do not have in the same font as the
answers I do. So this companion is organised as follows. Section~2 states the epistemic status of
every kind of claim in it. Section~3 reports four new computations. Section~4 grades the source pool by
kind, with the PBR verification as a worked example. Section~5 is the concept ledger: fifteen terms that
carry weight in the main paper, each defined exactly, then said again in plain words, then fenced off
from the misuse I have seen it suffer. Section~6 is the objection audit. Section~7 lists what I refuse
to conclude and what would change my mind, updated with the numbers below. Appendix~A is reproduction.

Two conventions from the main paper carry over unchanged, because they are the point of the exercise.
Every claim is tagged: \textbf{[FACT]} (established, checked against a primary source), \textbf{[FACT,
CITED, NOT FETCHED]} (established, but relying on my reading of a source I could not open in this
session), \textbf{[MY CALCULATION]} (arithmetic I ran here, with the log line findable by searching its
own digits), \textbf{[MY MEASUREMENT]} (a value obtained by fitting data I downloaded), and
\textbf{[SPECULATION]} (inference, mine or somebody else's, that outruns the above). And every
quantitative claim in this companion has a figure, five of them, with the caption generated from the
same JSON that the plot was drawn from.

\section{Evidence tiers, and what each one forbids}

A tier system is only worth having if it constrains the sentence that uses it, so I will state the
constraint rather than the definition. \textbf{[FACT]} may be used as a premise in an argument.
\textbf{[FACT, CITED, NOT FETCHED]} may be used to report a position but not to refute one --- I will
not tell a philosopher that they are wrong about their own paper if I have only read about it.
\textbf{[MY CALCULATION]} may be used to draw a conclusion only if the same paragraph states the model
that the calculation is internal to; my $10^{122}$ entropy comparison is a calculation inside
general relativity plus quantum field theory in a fixed background, and it buys no argument about
whether that is the right framework. \textbf{[MY MEASUREMENT]} comes with the additional requirement
that I report how many parameters I fitted and what the goodness-of-fit null looks like, which is
Section~3.1. \textbf{[SPECULATION]} may not appear in a sentence containing a number unless the number
is the thing being speculated about.

\section{Four new computations}

\subsection{Is my supernova fit ``too good''? A goodness-of-fit null, simulated}

The main paper reports a two-parameter fit to the 40 binned Pantheon distance moduli with $\chi^2 =
39.30$ on 38 degrees of freedom, and a reader is entitled to ask two opposite questions at once: is
that too high (is the model wrong) and is that too low (is my covariance overestimated, or did I
sneak in freedom)? A table of $\chi^2(38)$ answers the first question and an answer of my own making
answers the second.

I generated 4,000 synthetic 40-bin Hubble diagrams by taking my own best-fit model and adding
correlated noise drawn from the published $40\times40$ covariance, then refitting each one with
exactly the code that produced the number in the text. The null distribution of the refit statistic
has median $37.38$ and a 90\% band of $[25.18,\,53.04]$; my observed $39.30$ sits at the 58.6th
percentile, with $P(\chi^2_{\text{null}}\ge \chi^2_{\text{obs}}) = 0.414$ and $P(\le) = 0.586$, so the
two-sided ``too good'' $p$-value is $0.828$. \textbf{[MY CALCULATION]} The fit is unremarkable in both
directions (Figure~\ref{fig:gof}). As a cross-check on my own machinery, the textbook $\chi^2(38)$ tail gives $0.412$ against
my simulated $0.414$ --- agreement to $0.0021$, which tells me that profiling out the amplitude and
using a non-diagonal covariance do not shift the null in this case; the simulation is right by
construction and the table is what a referee will reach for, so I report both.

\begin{figure}[H]
\centering
\includegraphics[width=1.0\textwidth]{figures/f20_chisq_gof_null.png}
\caption{The goodness-of-fit null for my own fit, simulated rather than looked up. 4,000 synthetic
Hubble diagrams drawn from my best-fit model with the published covariance, each refitted by the same
routine that fitted the real data. Observed $\chi^2 = 39.30$ (red), null median $37.38$ (grey dashed),
90\% of null fits (orange). $P(\ge) = 0.414$, $P(\le) = 0.586$. \textbf{[MY CALCULATION]}}
\label{fig:gof}
\end{figure}

\subsection{How wide is my error bar, and does it depend on how I got it?}

The main paper quotes $\Omega_{\text{m}} = 0.2961 \pm 0.0220$ from a 300-draw parametric bootstrap. The
sceptic's question is whether that interval is an artefact of resampling. I computed it three more ways
on the same likelihood: a parabola fitted to the $\chi^2$ scan (curvature), the profile-likelihood
crossings of $\Delta\chi^2 = 1, 4, 9.21$, and the grid minimum itself.

The parabolic refinement of the grid minimum gives $\Omega_{\text{m}} = 0.29624$ with $\chi^2 = 39.2963$,
and a curvature-based $1\sigma$ of $\pm 0.02167$. The profile intervals are $1\sigma$:
$[0.2748,\,0.3182]$, $2\sigma$: $[0.2541,\,0.3411]$, $3\sigma$: $[0.2335,\,0.3655]$. The
half-width of the profile $1\sigma$ interval is $0.0217$ against the shipped bootstrap $0.0220$ ---
agreement to $1.3\%$ of the interval. \textbf{[MY CALCULATION]} So the acceleration claim in the main
paper (a matter density $16.8\sigma$ below the no-acceleration boundary $\Omega_{\text{m}}=2/3$) does not
rest on the resampling choice: profile, bootstrap and curvature give the same width, and a reader who
prefers a different convention can read the alternative straight off Figure~\ref{fig:prof}. The
$3\sigma$ interval still tops out at $0.3655$, which is a long way below $2/3$; no interval convention I
computed rescues a non-accelerating universe.

\begin{figure}[H]
\centering
\includegraphics[width=1.0\textwidth]{figures/f22_profile_likelihood_Om.png}
\caption{One curve, three confidence conventions, and the interval I originally shipped. $\chi^2$ of my
fit with the amplitude eliminated analytically, scanned over 313 values of $\Omega_{\text{m}}$; the
crossings of $\Delta\chi^2 = 1,4,9.21$ define the profile intervals, the dotted purple segment is the
300-draw bootstrap quoted in the main paper, and the dash-dot red line is the $\Omega_{\text{m}}=2/3$
no-acceleration boundary. \textbf{[MY CALCULATION]}}
\label{fig:prof}
\end{figure}

\subsection{What the pooling choice does to the number people fight about}

The main paper pooled measurements of $H_0$ and reported a random-effects mean of $70.561 \pm 1.304$
with $I^2 = 54.4\%$. The obvious objection is that any such number is an artefact of the estimator. I
answered it by running every defensible estimator on a pool I define explicitly here --- the four
strictly independent anchors: Planck 2018 $\Lambda$CDM $67.36\pm0.54$, SH0ES R22 $73.04\pm1.04$, the
CCHP TRGB distance $68.81\pm1.80$, and the JAGB $67.80$ with $2.17$ statistical and $1.64$ systematic
added in quadrature to $\pm2.720$ --- and by computing, rather than resampling, a Bayesian
random-effects model: flat prior on $\mu$ over a 401-point grid, half-Cauchy$(0,1)$ prior on
$\tau$ on a 401-point grid, the joint posterior evaluated on the full $401\times401$ lattice and
normalised, which is what ``marginalising over the heterogeneity'' should mean if it means anything.

Results \textbf{[MY CALCULATION]}: unweighted mean $69.2525$; fixed-effect inverse-variance mean
$68.5603\pm0.4565$ with $Q = 23.59/3$ and $I^2 = 87.3\%$; DerSimonian--Laird $69.3833\pm1.7372$ with
$\tau = 3.099$; Paule--Mandel $\tau = 1.527$; REML $\tau = 2.567$; Bayesian $\mu = 69.3870\pm1.6209$,
$95\%$ credible interval $[66.08,\,72.69]$, $\tau$ median $2.300$ with $95\%$ CrI $[0.96,\,5.94]$;
median of the four values $68.3050$. Four things follow, and they are the kind of thing the main paper
argues for but could not demonstrate in the space available.

\emph{One.} The four reasonable point estimators span $1.0820$~km/s/Mpc, which is $2.37$ fixed-effect
errors. Anyone quoting ``the'' pooled $H_0$ to three digits is quoting an estimator choice.
\emph{Two.} The Bayesian interval is $3.5$ times wider than the fixed-effect error, and the
five-estimator disagreement is entirely a disagreement about $\tau$ --- Figure~\ref{fig:lad}, right
panel, plots the three frequentist $\tau$ estimators against the whole posterior of $\tau$, which
stretches from $0.96$ to $5.94$. \emph{Three.} The posterior predictive interval for what the
\next\ measurement will say is $[62.51,\,76.29]$, a width of $13.8$~km/s/Mpc. That is the number to
quote when somebody says a new result ``resolved'' the tension: nothing about this pool makes a value
near $67.4$ or near $73.0$ surprising. \emph{Four.} Leave-one-out is brutal and I print it because it
is the test that should govern the whole genre. Dropping Planck moves the pooled Bayesian estimate to
$70.854$ ($+1.47$); dropping R22 moves it to $67.639$ ($-1.75$); dropping TRGB or JAGB moves it by
$+0.16$ and $+0.29$. The pooled centre can therefore be moved by $3.21$~km/s/Mpc --- $2.0$ posterior
standard deviations, and $4.08$ under fixed-effect weighting --- merely by deciding which experiments are in the room (Figure~\ref{fig:loo}).

\begin{figure}[H]
\centering
\includegraphics[width=1.0\textwidth]{figures/f19_pooling_estimtor_ladder.png}
\caption{Left: every defensible point estimator of the same four anchors, with the Planck and R22 input
values as horizontal references. Right: the three frequentist $\tau$ estimators against the full
marginal posterior of $\tau$ from the grid-integrated Bayesian model. The point estimators disagree by
$1.08$~km/s/Mpc ($2.37$ fixed-effect errors) and the entire disagreement is about this one nuisance
parameter. \textbf{[MY CALCULATION]}}
\label{fig:lad}
\end{figure}

\begin{figure}[H]
\centering
\includegraphics[width=1.0\textwidth]{figures/f21_leave_one_out_H0.png}
\caption{Leave-one-out. For each experiment in the pool, the shift in the pooled Bayesian $H_0$
relative to the measurement that was deleted. The pooled summary is not an independent witness: its
centre moves by up to $4.1$~km/s/Mpc depending on who is counted. \textbf{[MY CALCULATION]}}
\label{fig:loo}
\end{figure}

I must also state a discrepancy rather than bury it. This four-anchor pool gives \emph{larger}
heterogeneity ($Q=23.59/3$, $I^2=87.3\%$) than the number the main paper reported ($Q=6.58/3$,
$I^2=54.4\%$), because the main paper pooled seven items, including derived combinations that share
their inputs. Combining a measurement with a combination that contains it shrinks $Q$ by construction.
That is not a contradiction to smooth over; it is a further reason to distrust any pooled $H_0$,
including the one I shipped. I could not re-read the exact $(y,\sigma)$ vector behind the shipped
number in this pass --- the file that held it was lost with the workspace --- so the seven-item pool is
left as published in the main paper, and the four-anchor version above is the honest one. The same
reasoning applies to the $S_8$ pooling in section 3.7 of the main paper, and for the same reason I did
not re-run the $S_8$ ladder here: reconstructing that five-value vector from memory would let a
remembered input masquerade as a check, which is the failure mode this entire audit exists to avoid.

\subsection{How much of ``why now'' is my threshold, and what a reimplementation caught}

The main paper's ``why now'' paragraph computes the interval over which the matter and dark-energy
densities are within a factor of ten of each other, gets $18.21$~Gyr against a Hubble time of $14.53$
Gyr, and concludes that the coincidence is too broad to argue from. The question I did not answer
there: how much of that is my choice of ``a factor of ten''?

I recomputed the whole curve from scratch, in a different way. For flat $\Lambda$CDM the age is
analytic, $t(a) = \tfrac{2}{3}(H_0\sqrt{\Omega_\Lambda})^{-1}\,\mathrm{asinh}\big(\sqrt{\Omega_\Lambda
/\Omega_{\text{m}}}\,a^{3/2}\big)$, so both edges of the window follow from solving
$\Omega_{\text{m}}(1+z)^3/\Omega_\Lambda = R$ and $=1/R$. \textbf{[MY CALCULATION]} With $\Omega_{\text{m}} =
1-\Omega_\Lambda = 0.3153$: the ratio today is $0.4605$; the \emph{smallest} tolerance for which today
counts as coincidental at all is therefore a factor $2.1716$; and the windows are $18.214$~Gyr at
$R=10$ ($132\%$ of the present age), $8.978$~Gyr at $R=3$, $6.372$~Gyr at $R=2.171$, $5.704$~Gyr at
$R=2$, and $3.347$~Gyr at $R=1.5$ ($24\%$ of the age). Tightening the tolerance by a factor of six
shrinks the window by only a factor of $5.4$; and at $R=2$ the present day is \emph{outside} the
window, which is the sharpest possible statement of how soft this coincidence is. My $t(0)$ from the
closed form is $13.8011$~Gyr against the Planck quoted age of $13.7981$~Gyr, so the units and the
prefactor reproduce the standard result to $3$~Myr from Planck's parameters alone --- the check that
legitimises every Gyr in this subsection.

Now the confession, because it is the methodological payload. My first implementation of this
subsection used a quadrature of $\int \mathrm{d}z/((1+z)H)$ with the limits mirrored for negative $z$.
Its \emph{past} leg agreed with the main paper to $0.03$~Gyr, and its \emph{future} leg was wrong by a
factor of $1.8$ --- it returned $4.42$~Gyr for an interval of $8.05$~Gyr, and hence a window of $14.58$
Gyr instead of $18.21$. Had I kept the quadrature and only checked the past leg, as I nearly did, this
companion would have printed a number that contradicts the main paper, and the story would have been
about my error rather than about the tolerance curve. Two lessons generalise, and both are now stated
where they will be read: a code path that agrees on one side of a sign change proves nothing about the
other side; and an analytic form is worth a day of debugging precisely because it adjudicates. The
shipped $18.21$~Gyr survived, and Figure~\ref{fig:tol} plots the curve that the corrected code
produces.

\begin{figure}[H]
\centering
\includegraphics[width=1.0\textwidth]{figures/f23_whynow_tolerance_curve.png}
\caption{The ``why now'' window as a function of the tolerance, computed from the analytic age rather
than the quadrature used in the main paper, and agreeing with it to $0.004$~Gyr at $R=10$. Left: the
window in Gyr, with the Hubble time and the age of the Universe as reference lines, and the vertical
mark at $R=2.17$ where today ceases to qualify at all. Right: the cosmic age at the far edge of each
window as a fraction of today's. \textbf{[MY CALCULATION]}}
\label{fig:tol}
\end{figure}

\section{The source pool, graded by what kind of object it is}

A citation is not a unit of evidence; the \emph{kind of object} is. Six grades, with the instances from
the main paper that belong in each, and the constraint each grade imposes.

\textbf{(1) Datasets and their releases.} The Pantheon light-curve compilation and its
$1048\times1048$ systematic covariance, taken as a tarball from the maintainers' repository; the
supernova figures in the main paper and the fit interrogated in Section~3.1 above run on those files.
What a dataset release can bear: numbers you can recompute. What it cannot bear: the analysis. The
release I used deliberately has its light-curve shape and colour corrections \emph{already applied}
and those columns zeroed, which is why my raw 1,048-row fit and my 40-bin fit agree --- and why neither
can be extended into a measurement of the nuisance parameters.

\textbf{(2) Peer-reviewed primary results.} Planck~2018~VI, the R22 distance ladder, the COGITATE
adversarial trial in \emph{Nature} (2025), Albantakis et al.\ on IIT~4.0 in \emph{PLOS Computational
Biology} (19:e1011465, 2023), Tegmark in \emph{Phys.\ Rev.\ E} (61:4194, 2000), Casali et al.\ on the
perturbational complexity index. Grade 2 means the claims were read by referees, not that they are
right: the COGITATE result is grade~2 and it is also the example in the main paper of a flagship
adversarial test whose two camps each reanalysed the same data to opposite conclusions.

\textbf{(3) Peer-reviewed \emph{reanalyses} of other people's results.} The estimator-dependent
DCM reanalyses of COGITATE in \emph{eLife}, the Tremblay \& Lefebvre reanalysis of the
$\gamma$-ray/meditation anomalies, Salaberli \& Soto's appraisal of $\Phi$ in \emph{Neuroscience
Biobehavioral Reviews} (2025), the rotating-model lithium papers. This is the highest-value grade in a
contested area and the easiest to over-weight: a reanalysis inherits the primary data and the primary
model choices, and it can be wrong in a way that neither the original nor the critic noticed. When a
reanalysis and its target disagree, the main paper's rule is to report both and to refuse to adjudicate
unless I can compute something myself.

\textbf{(4) Self-reviewed preprints.} arXiv and bioRxiv items: the identifiability result against
$\Phi$, the consciousness-indicator programme, the 2026 DESI/YEAR-1 style drafts. A preprint is a
claim under construction, and the correct use is to cite what its own text says, never what a thread
about it says. Two concrete verifications, one of which is a lesson.

\emph{Worked case --- the PBR theorem.} In the last pass I drafted a sentence attributing to Pusey,
Barrett and Rudolph the structure ``two explicit premises, and a forced conclusion.'' I have now
checked the primary object directly on arXiv rather than trusting my recall of its identifier:
arXiv:1111.3328, submitted 14 November 2011, version~3 dated 18 November 2012, published as
\emph{Nature Physics} 8, 475 (2012) with DOI 10.1038/nphys2309, and the abstract states the conclusion
as ``any model in which a quantum state represents mere information about an underlying physical state
of the system, and in which systems that are prepared independently have independent physical states,
must make predictions which contradict those of quantum theory.'' \textbf{[FACT]} So the two premises
are exactly what the main paper says they are --- a $\psi$-epistemic reading of the state, and
preparation independence --- and the theorem is an implication between them and a contradiction with
quantum predictions, not a proof that the wavefunction is real in any metaphysical sense. The same API
query returned a second object worth reporting: a ``Comment on `On the reality of the quantum state'''
(arXiv:1812.03035, November 2018) claiming that PBR err at the start of their argument by violating
the projection postulate --- with the arXiv record annotated \emph{``withdrawn by the author; due to a
Babylonian confusion of tongues, the comment was based on an interpretation of statements in the PBR
paper that was not implied by its authors.''} \textbf{[FACT]} That is the whole argument for grade~4
discipline in one record: preprints carry the author's own retraction note, and a claim that survives
in citations after it has been withdrawn by its author is exactly how a literature acquires a
disagreement it never had.

\emph{Counter-case --- a citation I refuse to use.} In an earlier draft I had a line attributed to a
\emph{Nature} volume and page with a Bayes factor of $44.96$. I could not locate the object at that
address in this session. The main paper's ledger does not cite it and neither does this one; a
number whose container I cannot open is not evidence, however much it supports me.

\textbf{(5) Curated reference works.} The \emph{Stanford Encyclopedia of Philosophy} entries (for this
pass, the panpsychism article, first published 2001, substantive revision 13 May 2022), which is where
I took the contemporary taxonomy used in Section~5: the constitutive/emergentist split, panpsychism
versus panprotopsychism, micropsychism versus cosmopsychism, Russellian monism, the anti-emergence and
intrinsic-nature arguments, and the three standard objections --- the incredulous stare, the combination
problem, and the subject-summing problem. A reference work is not a primary argument and cannot be cited
for one; it is a check that I have not invented a taxonomy or forgotten a standard objection. Where the
main paper says the combination problem is ``the hardest problem in the view,'' the encyclopedia's own
framing is what I am relying on, and the primary literature (James for the problem, Goff and Seager for
contemporary attempts) is cited at grade~5 status: read as positions, not used to refute.

\textbf{(6) Everything else --- interviews, journals of speculative metaphysics, secondary books,
press material.} Used in the main paper exactly once, to characterise a genre's rhetoric rather than
its content. If a claim in either document is graded~6, it is decoration and I have said so in place.

\section{The concept ledger: fifteen terms, each defined twice and fenced once}

The main paper's sin would be to let these terms slide between registers. For each: the exact sense,
the plain-words sense, and the misuse to be avoided.

\textbf{Distance modulus} $\mu = m - M = 5\log_{10}(d_L/10\,\mathrm{pc})$. In words: how many
magnitudes fainter a standard candle looks than it would at 10 parsecs, which is a logarithm of
distance. Not a distance: a $\mu$ without a stated $H_0$ convention fixes the \emph{shape} of the
expansion history and nothing about its scale, which is the whole content of the main paper's
supernova section.

\textbf{Sound horizon} $r_{\text{d}}$: the comoving distance a pressure wave could travel in the photon
baryon fluid before the plasma decouples at recombination, i.e.\ $\int_{z_{\text{d}}}^\infty c_s(z)\,
\mathrm{d}z/(\dot a/a)$, with $c_s$ dragged down by the baryon loading. In words: the length of the
ruler printed in the CMB, set by how far sound could travel before the Universe turned transparent.
Not a standard rod of known length independent of cosmology: it is $r_{\text{d}} = r_{\text{d}}(H_0\sqrt{
\Omega_{\text{m}}},\ldots)$, so the ``inverse'' $H_0$ from BAO inherits a model, and that lock is what
Figure~4 of the main paper draws.

\textbf{Degeneracy versus uncertainty.} A degeneracy is a direction in parameter space along which the
likelihood is nearly flat --- the amplitude $M$ and $H_0$ in my fit; the ``curvature degeneracy''
$-\Omega_k$ traded against $H_0$. An uncertainty is a width along a direction that does exist.
Reporting a degeneracy as an error bar manufactures a crisis out of a ridge.

\textbf{$Q$, $I^2$, $\tau$.} For a set of measurements: $Q$ is the weighted scatter about the
fixed-effect mean, compared to its null; $I^2$ says what fraction of the observed scatter is
real between-study variation rather than error bars; $\tau$ is that extra scatter, in the units of the
quantity. In words: $\tau$ is the amount by which the experiments disagree with each other beyond what
their own quoted errors allow --- $2.30$~km/s/Mpc median for my $H_0$ pool, and the quantity every
pooled number hides.

\textbf{Profile likelihood versus bootstrap.} The profile interval is the set of parameter values that
survive when the nuisance parameters are re-optimised at each value; a bootstrap resamples the data.
The first is geometry of the likelihood surface, the second is a statement about the sampling
procedure. Section~3.2 shows they agree here; where they disagree, the disagreement is the finding.

\textbf{Posterior predictive interval.} The interval for a \emph{new} measurement, integrating over
the parameter uncertainty and the between-study scatter. In words: given what we know, what should the
next experiment find? If the answer includes both camps' numbers, the pool has not resolved anything,
which is what mine returns: $[62.5,\,76.3]$.

\textbf{Decoherence.} Suppression of off-diagonal density-matrix elements in a preferred basis through
entanglement with uncontrolled environmental degrees of freedom, at a rate set by the system-environment
coupling. In words: information about a superposition leaks into the room faster than you can follow
it, so the interference becomes unavailable in practice. Not a collapse: no single outcome is produced
by decoherence, and no observer is required by it --- the most common conflation in the quantum-mind
literature, and the one the main paper's arithmetic section exists to keep separate.

\textbf{Reduced state / partial trace.} The state assigned to a subsystem after averaging over its
environment; $\rho_S = \mathrm{Tr}_E \rho_{SE}$. In words: what the part looks like when you are not
allowed to look at the rest. It is how ``the brain is warm and wet'' becomes a number, and the number is
$Q=\nu\tau$, not a mood.

\textbf{Measurement problem versus collapse.} The problem is that the theory has two evolution laws and
no stated criterion for which applies. A collapse model is a specific \emph{solution}: a modified
dynamics with a parameter, hence a prediction, hence a way to lose. ``Consciousness causes collapse''
is not a member of the second class: it changes nothing you can count.

\textbf{The hard problem.} The question of why there is something it is like to be a physical system
that processes information, and why that ``like'' has the properties it does. In words: even after you
describe the mechanism of colour vision completely, it remains sensible to ask why seeing red feels
like something. \emph{Not} the hard problem: the explanatory gap as a feeling of confusion, and the
access problems (report, attention, integration) that are empirically tractable --- Chalmers' own
distinction, and the reason COGITATE could only test correlates.

\textbf{Epistemic versus metaphysical possibility.} ``Conceivable'' can mean ``I can construct a model
of it'' or ``it could be true.'' The conceivability arguments for dualism and against physicalism need
the second reading to work and only ever establish the first. Zombies and Mary are arguments, not
intuitions, and they have a specific target: the inference from the complete physical to the
phenomenal.

\textbf{$\Phi$ versus complexity versus PCI.} $\Phi$ (IIT's integrated information) is a property of a
mechanism's cause-effect structure under a maximum over subsystem bipartitions, computed on a
transition probability matrix. Lempel--Ziv complexity counts the phrases needed to reconstruct a
bit-string. PCI is a normalised Lempel--Ziv complexity of a binary perturbation-response matrix. In
words: $\Phi$ claims to measure how much a system is \emph{one} thing; the other two measure how
\emph{variegated} a recorded signal is. They coincide on nothing in my toy computations, and the main
paper times the difference: $\Phi$ is $O(2^n)$ with a measured $\ln t = 3.05 n - 9.92$, while
Lempel--Ziv is linearithmic.

\textbf{Combination versus dissociation.} Micro-experientialism must build one subject from many;
cosmopsychism must divide one subject into many. In words: the first owes an account of how
experiential grains merge, the second owes an account of why the boundaries fall where skulls do.
Both are real problems, and neither has a working solution, which is why the main paper calls this the
load-bearing distinction rather than a matter of taste.

\textbf{Intrinsic versus relational property.} A property a thing has in itself versus one it has
only relative to something else; QBism reads the quantum state as the second kind. In words: whether
the wavefunction is the cat's or the bettor's. The measurement problem's shape changes with this
choice, and no experiment I know of distinguishes them --- which is precisely the sentence the
main paper's verdict rests on.

\textbf{de Sitter entropy versus thermodynamic entropy.} $S = A/4\ell_{\text{P}}^2$ for a cosmological
horizon: a counting of horizon degrees of freedom, $3.3\times10^{122}$ in $k_B$ units in my
arithmetic. In words: an upper bound on what an observer inside the horizon can know or store. Not the
entropy of the contents of the Universe, and not a licence to compare it with $10^{88}$ as if the two
were counts of the same kind --- which is exactly why the Boltzmann-observer argument in the main paper
is presented as a budget constraint and not as a refutation.

\section{The objection audit}

Each entry: the question as a reader would pose it, what an honest answer would require, and the
answer. Where I have no answer, that is the answer.

\textbf{1. ``You fitted 40 bins with 2 parameters and got $\chi^2/{\text{dof}}\approx1$; of course you did --- you chose the model that the bins encode.''} Correct, and it is a limitation of every
supernova analysis including the collaboration's: Pantheon's bins are produced under an assumed
flat-$w$CDM family. What my fit can legitimately support is a \emph{consistency} claim and the
acceleration boundary, not an independent test of dark energy. Stated in the main paper, restated here.

\textbf{2. ``Your covariance is a published approximation, not a likelihood.''} Yes. The binned
systematic matrices are given as $N\times N$ covariance blocks over one parameter family, and for the
full 1,048-object release my own eigenvalue check finds a minimum eigenvalue of $+7.4\times10^{-3}$,
so no regularisation was needed there, but the 40-bin matrix is not the covariance of the bin means
under the full analysis. My numbers are thus $\chi^2$ values under an approximation, which is why I
report $1\sigma$ to three decimal places and never a $5\sigma$ claim.

\textbf{3. ``Why not just use Pantheon+ with the bias corrections?''} The file I could retrieve from the
public repository carries a \texttt{biascor} column of zeros; the Pantheon+ analysis requires the
host-mass and calibration nuisance parameters whose values are not in the text file. Using the
corrections I cannot read would make my numbers somebody else's, which is the failure mode I am
auditing for.

\textbf{4. ``You quote the tension at $4.8\sigma$ from Planck versus R22 but pool to $2.3\sigma$;
you are averaging away the physics.''} That is precisely my claim, and Section~3.3 sharpens it with the
estimator ladder: the pooled value depends on the estimator by $2.37$ fixed-effect errors and on
membership by $4.1$~km/s/Mpc. A tension reduced by pooling is not resolved, it is averaged, and an
average of two clusters of measurements with $\tau=2.3$ is not a prediction.

\textbf{5. ``Your Bayesian model's half-Cauchy prior on $\tau$ is a choice that inflates the interval.''}
It is a choice, and I state it: Gelman's recommendation for scale parameters, weakly informative, and
the alternative estimators of $\tau$ (DL, PM, REML) span $1.53$ to $3.10$, which brackets the posterior
median of $2.30$. If a reader prefers a different prior, Figure~\ref{fig:lad} right panel gives the
whole marginal so they can recompute the interval by eye.

\textbf{6. ``Why does your Hubble-time arithmetic differ from the main paper's 14.53 Gyr?''} Because I
used $\Omega_{\text{m}}=1-\Omega_\Lambda=0.3153$ from Planck's $\Omega_\Lambda$ rather than the paired
$\Omega_{\text{m}}=0.3142$; the resulting difference is $0.014$~Gyr, $0.1\%$, and it is reported rather
than smoothed, as is the $0.004$~Gyr window difference in Section~3.4.

\textbf{7. ``The lithium problem is the strongest crack you cite and you refuse to call it a crisis.''}
Yes: the predicted value is primordial, the observed value is a stellar surface abundance, and
2025--26 stellar-model and chemical-evolution work argues the discrepancy may be an artefact of which
number gets compared to which. A $6\sigma$ offset between two quantities that are not the same
quantity is not a crisis in $\Lambda$CDM. That is a fact about the objects being compared, not a
judgement about who is right.

\textbf{8. ``DESI's dynamical dark energy: your paper is out of date.''} The main paper reports the
Year-1 result and the later data release's changed preference, and its programme section commits to
checking the next release. What I will not do is quote a tension in $w$ at a significance I cannot
reproduce from a covariance I do not have.

\textbf{9. ``Page and Geilker refuted Wigner, so why spend a section on it?''} Because the experiment
is the only case where a philosophical distinction --- the cut's placement --- was made into a
measurement, and it came back against the mind-causes-collapse reading. The scope limit is stated in
the main paper: it constrains collapse triggered by conscious observation as the completion of a
measurement; it does not touch Everett, Bohm, QBism, or any interpretation that makes no such dynamical
claim.

\textbf{10. ``Bong et al.\ 2020 shows local friendliness fails --- doesn't that bring the observer in
after all?''} It closes a trichotomy (locality, freedom of choice, and observer-independent facts) under
assumptions including no superdeterminism, with a laboratory implementation. \textbf{[FACT, CITED, NOT
FETHED]} I did not re-verify its exact premise list in this pass --- and that is the honest status
marker for the claim in the main paper, which uses it to say that the question is now a theorem-shaped
question rather than an intuition.

\textbf{11. ``Is the decoherence objection to quantum consciousness dead?''} Tegmark's $10^{-13}$~s
estimate was shown to rest on a dipole moment roughly two orders of magnitude too small; Hagan et al.\
raise the timescale toward $10^{-4}$ to $10^{-5}$~s. My contribution is arithmetic, not physics: at
$\nu\tau$ with $T=310$~K, a 613~THz mode is $95\,k_BT$ and thermally unpopulated by some 42 orders of
magnitude, so a Fr\"ohlich-type coherence needs a pump; and $1612$ dimers relaxing per microtubule
against a $4.74\times10^{-12}$~W-per-neuron budget is the kind of sentence that must be checked, not
recited. The objection is not answered by the corrections --- the mode is still unpopulated --- and the
correction is not answered by the objection either. Both camps skip the arithmetic.

\textbf{12. ``You timed $\Phi$ on toys; what does that prove?''} Nothing about IIT's truth. It
constrains its use: a measure whose exact evaluation on $n$ elements is exponential in $n$ (my
measured $\ln t = 3.05n - 9.92$, $t$ at $n=8$ about three weeks, $n=10$ about 28 years) cannot be the
operative quantity behind a claim about the cortex, and cannot be the basis of an ethical weight for
AI systems. Combined with the identifiability preprint and the Salaberli \& Soto appraisal, that is the
main paper's verdict on the measure while remaining agnostic about the axioms.

\textbf{13. ``If IIT is unfalsifiable, why treat it with so much space?''} Because its axioms are
defensible phenomenology and its postulates are a non-trivial bridge, and it is the only theory in the
field that has published the bridge's cost. The falsifiable version of the analogy in the main paper
(the $\Phi$/PCI dissociation) is a prediction about measurement, and if the next generation of
measures kills it, the section will need rewriting.

\textbf{14. ``Your PCI reimplementation gives complex/white-noise ratios of 1.008 --- are you saying
the published index is broken?''} I am saying that \emph{the estimator as I reimplemented it from the
published description} does not resolve the contrast its own figure advertises, that it drifts by
$0.117$ versus $0.112$ between two noise draws, and that two independent draws of the same generator
differ by $0.0565$. \textbf{[MY CALCULATION]} Whether the published implementations avoid this is not
something I can answer: I did not have their code. That is a limitation of the objection, not of the
finding.

\textbf{15. ``Clinicians use PCI with a threshold around 0.4; are you dismissing that?''} No: the main
paper's strongest consciousness claim is precisely the validated one --- Casali et al.'s
$0.45$--$0.67$ versus $0.25$--$0.31$ across consciousness states, with a criticality-flavoured
cut-off near $0.35$ in the later LOSO analysis. What I dismiss is the further step from ``this
measure tracks a state'' to ``this measure quantifies experience.''

\textbf{16. ``Why no free will, no ethics, no AI consciousness section?''} Because the brief was the
Universe and consciousness, and each of those three would require evidence I do not have. The main
paper mentions the indicator programme (Butlin et al.) only to say the method is a checklist of
mechanisms, not a measurement, and that is the whole extent of what I can support. Panpsychism gets
space because it is a live position in the philosophy of mind with an explicit combination problem; a
theory of machine consciousness does not, in my reading of the literature.

\textbf{17. ``You say the philosophy is [CITED, NOT FETCHED] but you cite forty years of it. Where is
the risk?''} Everywhere, and it is stated per sentence. The concrete risk I identified while preparing
this companion is arXiv-identifier drift: three identifiers I wrote from memory resolved on the arXiv
API to a HST detector-calibration paper, an image-captioning paper and a post-5G networking paper.
Nothing in either document depends on those identifiers; the PBR entry in Section~4 is the corrected,
verified form. If any arXiv number in the main paper's ledger conflicts with Section~4 here, Section~4
wins, because it is the one I opened.

\textbf{18. ``What would make you abandon the thesis of the main paper?''} One of three things: a
consciousness measure whose value is invariant under reanalysis with different model assumptions and
which dissociates from $\Phi$ in the direction I predict; a decoherence calculation for a plausible
neuronal substrate with $Q>1$ and an energy budget that closes; or a $1/H_0$ measurement from a
standard-siren or BAO-with-independent-$r_{\text{d}}$ route that lands near $67$ with an error bar under
$0.7$ and a documented refusal to use the Cepheid rung. Each is stated with the number attached in the
main paper's commitments section, and each is the kind of thing that could plausibly happen by 2031.

\section{What I still refuse to conclude}

I will not say the Universe is fine-tuned for observers; the tolerance curve in Section~3.4 is the
reason --- the coincidence survives at the price of any tolerance above a factor of $2.17$, which is
not a constraint on physics. I will not say the vacuum-energy discrepancy is $120$ orders of a
prediction, because it is a comparison between a low-energy field-theory estimate and a geometric
quantity, with the slope-4 comparison in the main paper showing how the number moves with an
ultraviolet cut-off I cannot justify. I will not say consciousness is measured, because the one
instrument with clinical validity measures a signal's variegation, not a subject's interior. I will not
say the $H_0$ tension is resolved or real, because I have now shown how much of it is my estimator: the
same four numbers support $68.31$ and $69.39$ and a $95\%$ credible interval that covers both camps'
headline values.

What the pair does support is weaker and, I think, more durable: the Universe is comprehensible in the
narrow sense that its expansion history can be recovered by two parameters fitted to 40 binned means by
a person with a laptop, and consciousness is not comprehensible in that sense yet, and the gap between
those two facts --- not a theorem, not a coincidence, not a $10^{122}$ entropy ledger --- is the actual
subject of the field. Everything else I have written in both documents is bookkeeping around that
sentence.

\appendix
\section{Reproduction}

The five figures and every number above come from one script:

\begin{verbatim}
cd research
python deepen.py 2>&1 | tee logs/deepen_log.txt   # ~3 s; writes paper/figures/f19..f23
                                                  # and logs/deepen.json (all numbers in machine form)
python paper/build.py                             # tex -> typst -> pdf; checks for markup leaks and unresolved refs
\end{verbatim}

\texttt{deepen.py} is self-contained: it re-reads \texttt{data/lcparam\_DS17f.txt} and
\texttt{data/sys\_DS17f.txt} (retrieved in this pass from \texttt{codeload.github.com/dscolnic/Pantheon},
the \texttt{Binned\_data/} directory of the master branch), re-implements the luminosity distance by
trapezoid integration of $1/E(z)$ with the $(1+z)$ factor and the radiation term, refits $\Omega_{\text{m}}$
over a 313-point grid with the amplitude eliminated analytically, computes the profile-likelihood
crossings with \texttt{brentq}, runs 4,000 parametric-bootstrap replicates of the fit statistic, solves
the Paule--Mandel and REML equations for $\tau$, integrates the Bayesian posterior on a
$401\times401$ grid, and computes the ``why now'' windows from the closed-form age $t(a)$. The caption
text of each figure is written to \texttt{logs/deepen\_figures\_manifest.json} by the same code that
draws the plot, and \texttt{logs/deepen.json} holds every number in machine-readable form, so a
disagreement between text and plot is detectable by diff rather than by argument.

Three checks a reader can perform in one line each. My grid fit should return
$\Omega_{\text{m}}=0.2950$ at $\chi^2=39.300$; the parabolic refinement should return $0.29624$; and
\texttt{t\_age\_gyr(0.0)} should return $13.8011$~Gyr, which is the Planck age reproduced from
Planck's $\Omega_\Lambda$ and $H_0$ alone. If those three agree, the numbers in this companion are
the numbers in its figures; if they do not, my Python is not doing what my prose says, and the
prose loses.

\vspace{6pt}
\hrule
\vspace{6pt}
\emph{Status of this document.} Section~3 is computation. Section~4 grades sources, with one
verification done in-session (PBR, via the arXiv API, plus the withdrawn-comment record) and the rest
carried at the grades stated. Section~5 and Section~6 are the reader's contract: definitions fenced
against misuse, and objections answered or marked unanswerable. No number in either document depends
on a philosophical source; no philosophical claim in either document is supported by a number.
\end{document}
